\documentclass[12pt]{article}
\usepackage[utf8]{inputenc}
\usepackage[T1]{fontenc}
\usepackage[margin=1in]{geometry}
\usepackage{setspace}
\usepackage{xcolor}
\usepackage[colorlinks=true, urlcolor=blue!60!black, linkcolor=blue!60!black]{hyperref}
\setlength{\parindent}{0em}
\setlength{\parskip}{0.45em}
\renewcommand{\baselinestretch}{1.0}
\begin{document}
\begin{center}
{\Large \textit{Gundy v. United States}, 588 U.S. 128 (2019)\par}
\vspace{0.3em}
{\normalsize Case Brief --- One-Page Edition\par}
\end{center}

\textbf{Memory Jogger.} SORNA's delegation to the Attorney General over pre-Act sex offenders survived because it meant registration ``as soon as feasible'' (\href{https://supreme.justia.com/cases/federal/us/588/17-6086/}{Justia}).

\textbf{Facts.} SORNA (2006) created a national sex-offender registry and let the Attorney General ``specify the applicability'' of its requirements to some 500,000 offenders convicted before the Act; he applied it to all. Herman Gundy, convicted in 2005, moved to New York after release without registering.

\textbf{Procedural History.} Convicted of failing to register, Gundy lost his nondelegation challenge in the district court and Second Circuit.

\textbf{Votes.} Affirmed, 5--3: Kagan for four; Alito concurred in the judgment; Gorsuch, with Roberts and Thomas, dissented; Kavanaugh did not participate.

\textbf{Holding.} Section 20913(d) is constitutional: construed to require registering all pre-Act offenders as soon as feasible, it supplies an intelligible principle.

\textbf{Reasoning.} The Act's comprehensive purpose, definitions, history, and reference to offenders unable to comply confine the discretion to timing; so read it is ``distinctly small-bore'' beside ``public interest'' standards long upheld. Alito would reconsider the doctrine with a willing majority. Gorsuch would allow only filling up details, fact-finding, or delegations within another branch's own powers, and accused the plurality of rewriting the text. \emph{Later:} \textit{FCC v. Consumers' Research} (2025) reaffirmed the intelligible-principle test, 6--3.

\textbf{Examples.} \emph{Same-Side:} a transitional drone-registry delegation. \emph{Opposite-Side:} sole discretion over whether a felony applies. \emph{Fence-Sitter:} AI felonies defined by a Secretary in the public interest.

\textbf{Critique.} Saving constructions may rewrite statutes, and a ten-year felony got no more scrutiny than utility ratemaking; unable to draw a nondelegation line, the Court turned to removal power and the major questions doctrine.

\textbf{Key Quotation.} ``Given that standard, a nondelegation inquiry always begins (and often almost ends) with statutory interpretation. Only after a court has determined a challenged statute's meaning can it decide whether the law sufficiently guides executive discretion to accord with Article I.'' --- Justice Kagan, plurality opinion, 588 U.S. at 135--36.

\end{document}
